\thispagestyle{empty}
\begin{center}{\bfseries\Large ABSTRACT\par}\end{center}
\vspace{4pt}
\begin{spacing}{1.25}\small
When a microservice fails, the failure does not stay where it starts. A slow datastore drags down the services that query it, their callers time out, retries amplify the load, and within a minute the operator faces dozens of alarms in which the originating fault cannot be told from its victims by severity alone. Recent diagnosis methods localise the fault with graph neural networks over fused metrics, traces and logs; they need labelled failures, GPU inference, and give an output the on-call engineer cannot audit.

This project builds the alternative as a complete, unsupervised, single-core pipeline that observes Docker containers from outside the application, over SSH, and assists the operator with a cause, a prediction of what fails next, a fault type and a runbook. A platform of about 36\,000 lines of TypeScript and 329 automated tests recovers the call graph from container socket tables, shows it in 2D and 3D, and runs the diagnosis. Its core is a probabilistic cascading-failure model whose per-edge propagation probabilities are learned online as Beta--Binomial posteriors; a hybrid detector in which a per-service Isolation Forest and a robust rule layer jointly raise the alarm, with container facts (stopped, crash-looping) treated as evidence rather than degree; a ranker that separates a service's own deviation from what its neighbours explain; and call directions learned from telemetry. A Detection engines page runs the Rules, ML and Hybrid engines in parallel and shows what each detected, computed and predicted.

The evaluation uses 62 faults injected blind into two real containerised systems over 217 minutes each, and 178 held-out synthetic faults over 16 topologies, with bootstrap intervals, paired tests, an ablation and an amplification study. The hybrid names the true cause first for 74\% of real faults against 65\% for the rule layer (paired difference $+9.7$ points, 95\% CI $[+1.6,+17.7]$); learned propagation discriminates better than fixed edge probabilities (AUROC 0.847 against 0.784) on 80 observed cascades; learned directions stay correct when service names carry no information (28 of 31 pairs against none); and change-point alarms roughly halve the median detection delay (10.5\,s to 4.9\,s). A decision window costs about 6 to 7\,ms. The report also gives what did not work, a measurement bug that had hidden half of the real faults, and operational findings from the benchmark hosts.
\end{spacing}
\clearpage
